\chapter*{前言}
\addcontentsline{toc}{chapter}{前言}

这本书不按“数组、链表、动态规划、系统设计”来编排，而按多邻国真实面试轮次来编排。原因很简单：多邻国的 Android 面试不是一场换皮的 FAANG 算法考试。你当然仍然需要会写代码、会分析复杂度、会在共享编辑器里把一个函数做对；但真正拉开差距的，往往不是那道中等难度题，而是候选人是否知道自己即将进入怎样的一天。

多邻国的循环有两个决定性的不同。第一，它有一轮独立的 Code Review，也就是很多中文候选人口中的“查错”。这不是让你从零写一个算法，而是让你读一段别人已经写好的代码，找出隐藏 bug、可维护性问题、边界条件、异步状态和 Kotlin 惯用性问题，然后像真实团队评审一样说清楚“为什么这是问题、应该怎么改、风险在哪里”。第二，对 Android 岗位来说，Design Interview 不是后端分布式系统设计，不是让你画一个全球 feed、支付系统或短链服务；它更接近 Android 架构设计：模块边界、状态流、离线与同步、Repository、ViewModel、依赖注入、Server-Driven UI、性能与可测试性。这两轮都很像真实工作，也都很不像 LeetCode。

这正是很多候选人的准备错位之处。大家把最多时间花在算法题库上，因为题库清晰、反馈快、刷题有成就感；却把 Code Review 和 Android 架构设计留到最后，用几篇零散面经临时补。结果到了现场，才发现“查错”轮要求的是在陌生代码里快速定位与表达，设计轮要求的是把移动端工程复杂度讲清楚，而不是背几个缓存和消息队列名词。本书因此把 17 道调试/查错题和 9 道 Android 设计题放在中心位置。算法仍然会讲，但它不是这本书的全部；更准确地说，算法只是你进入面试的门票，Code Review 和 Design 才是多邻国 Android 面试最值得提前投资的地方。

\begin{idea}[讲怎么想，而不是背答案]
这本书的目标不是给你一组“标准答案”。面试官真正想知道的是：你面对不完整需求时怎样澄清范围；面对陌生代码时怎样建立假设；发现 bug 后怎样判断严重性；提出设计后怎样承认取舍。你要练的不是把答案背熟，而是把推理路径说得让别人愿意相信。
\end{idea}

本书会非常强调来源可靠性。凡是来自多邻国官方博客、官方招聘说明或官方工程文章的信息，我会标成 \srcofficial；来自公开面经、候选人复盘、面试平台汇总的信息，会标成 \srccommunity；根据 Android 行业常识、多邻国公开技术栈和题型逻辑做出的判断，会标成 \srcinferred。请把这三类信息分开看。任何一本面经书如果不区分“官方确认”“候选人报告”和“作者推断”，都是危险的：它会让你把传闻当事实，把某一年某个人遇到的题当成固定流程，甚至在行为面里引用并不存在的公司价值观。

尤其要提醒你：多邻国官方工程博客不是背景读物，而几乎就是题库。Android Reboot、Kotlin 迁移、Server-Driven UI、Frontend Prediction 这些文章，直接暴露了他们关心的工程问题：全局可变状态、主线程阻塞、模块化、Kotlin 风格、服务端驱动 UI、乐观更新、一致性和公平性。后面的章节会反复回到这些材料。读这本书之前，至少先把第零章列出的三篇核心官方工程文章读一遍；读完后你会发现，很多所谓“面经题”并不是凭空出现的，它们只是这些真实工程经历在面试里的压缩版本。

最后，把这本书当成一本训练“面试现场表达”的书，而不是一本速查手册。每一轮你都要练三件事：先说清楚目标，再说清楚约束，最后说清楚取舍。如果你能让面试官感觉“这个人今天就在我的代码库里也能推进事情”，那才是多邻国 Android 面试真正想筛出来的信号。